% Regression: a leg the picture's physical policy grows is a named wire
% record, so it can be spoken about.  The policy used to write a field on
% each frame-cell atom and let the renderer ink a path nothing could name,
% while the same leg written as a typed port carried both a label and an
% address (issue 5574).  The leg now owns a record under the canonical name
% the crossing grammar already spells -- leg-<face>-<row>-<col> -- and every
% consumer of a name reaches it: a label and a bead take a station on it, and
% a declared crossing names it as an operand.
\documentclass{standalone}
\usepackage{tenkz}
\begin{document}
\tenkzkernel
% The policy's own legs, north and south of a two-site row.  Neither is
% deferred by a crossing, so each is inked by the leg pass and answers along
% the route it inked.
\begin{tenkz}[rows={wire}, cols=2, physical=updown]
  \tn{A} & \tn{B}
  \tnmark[form=label, label pos=0]{on leg-n-1-1 0.5}{$\sigma$}
  \tn[skin=dot, at=on leg-s-1-2 0.6]{}
\end{tenkz}
\qquad
% The same leg named by a declared crossing, which defers its ink to the
% string surgery.  The record answers a place on it either way.
\begin{tenkz}[rows={wire}, cols=3, bonds=none, physical=down]
  \tn[at=(1,1), name=Q, ports={0:virtual}]{Q}
  \tn[at=(1,2), name=M]{M}
  \tn[at=(1,3), name=L]{L}
  \tnwire[
    kind=string,
    name=g,
    route={s of (1,1) .. (1,3)},
    crossing=over
  ]{open}{open}
  \tnwire[
    kind=string,
    name=h,
    route={all of Q},
    crossing=over,
    cross={
      under at crossing of g and h,
      over at crossing of h and leg-s-1-2,
      over at crossing of h and leg-s-1-3
    }
  ]
    {crossing of g and leg-s-1-3}{1 n of Q}
  \tnmark[form=label, label pos=180]{on leg-s-1-2 0.3}{$\tau$}
\end{tenkz}
\end{document}
